\section{Existence of RARE}
\label{sec:exist}

In this section, we prove the existence of (approximate) RARE.
%
We first show that when we fix an allocation rule that is \emph{continuous},
then a $\kappa$-RARE always exists, where $\kappa$ is 1 for the rivalrous setting
and $b_{\avg}/b_{\min}$ for the non-rivalrous setting.

Then, we ask \emph{which} allocation rules are continuous,
and that leads us to different answers
for the rivalrous and non-rivalrous settings:
\begin{enumerate}
\item \textbf{Non-rivalrous setting}:
    When buyers' utility functions are strictly concave, the demand set is a singleton at all prices,
    so there is a unique allocation rule. Moreover, this allocation rule is continuous.
    Thus, $(b_{\avg}/b_{\min})$-RARE always exists for strictly concave utilities.
\item \textbf{Rivalrous setting}:
    Even with strictly concave utilities, the demand set may not be a singleton.
    But we identify a large class of natural allocation rules that are continuous,
    proving that RARE exists for all of them.
\end{enumerate}

For linear utilities, we present instances where no allocation rule is continuous.
Nevertheless, we prove that for concave (but not necessarily strictly concave) utilities,
a variant of RARE holds.

\subsection{Refining the Goal}

A price vector $\pvec \in \mathbb{R}_{\ge 0}^m$ is $\kappa$-RARE iff
\begin{quote}
for every seller $j \in [m]$ and deviation $q_j \ge 0$,
\\ there exists an individually rational response
$\qvec_{-j} \in \mathbb{R}_{\ge 0}^{[m] \setminus \{j\}}$, such that
\\ $R_j(\qvec) \le \kappa R_j(\pvec)$.
\end{quote}

We now sketch how to make the problem simpler by restricting our search space.
First, we restrict our search for a $\kappa$-RARE price vector $\pvec$ to the simplex $\beta\Delta_m$.
This puts it at the intersection of the low-price and high-price regimes from \cref{thm:accounting}:
prices are low enough that buyers can afford all goods,
but high enough to extract a $1/\kappa$ fraction of the total budget $B$.

Next, we show how to restrict the deviation $q_j$ and the responses $\qvec_{-j}$.
Fix an allocation rule and pick any price vector $\pvec \in \beta\Delta_m$.
By \cref{thm:accounting}.\ref{item:rev-props:low-prices},
each seller $k \in [m]$ earns revenue $R_k(\pvec) = \nu p_k$.
If $q_j < p_j$, then the deviator $j$'s revenue will drop from $\nu p_j$ to at most $\nu q_j$,
regardless of what the response $\qvec_{-j}$ is.
Thus, we need $q_j \ge p_j$ for the deviation to make sense.

For every respondent $k \neq j$, her revenue immediately after $j$'s deviation is
$R_k((q_j, \pvec_{-j})) \le \nu p_k$.
If every respondent $k \neq j$ raises her price to $q_k \ge p_k$,
the total revenue rises (or remains the same), i.e., $R_{\all}(\qvec) \ge R_{\all}(\pvec)$.
Then a response that preserves every seller $k$'s market share $R_k(\cdot)/R_{\all}(\cdot)$
would be individually rational, since
\[ R_k(\qvec) = \frac{R_k(\pvec)}{R_{\all}(\pvec)} R_{\all}(\qvec)
    \ge R_k(\pvec) = \nu p_k \ge R_k((q_j, \pvec_{-j})). \]
Moreover, we have $R_{\all}(\pvec) = \nu\beta$ and $R_{\all}(\qvec) \le B$, giving us
\[ R_j(\qvec) = R_j(\pvec) \frac{R_{\all}(\qvec)}{R_{\all}(\pvec)} \le \kappa R_j(\pvec). \]
The original market share of each seller $k \in [m]$ is $R_k(\pvec)/R_{\all}(\pvec) = p_k/\beta$.
Hence, for $\pvec \in \beta\Delta_m$ to be a $\kappa$-RARE, it suffices to show that
\begin{quote}
for every seller $j \in [m]$ and deviation $q_j \in [p_j, \infty)$,
\\ there exists a response $\qvec_{-j} \in \prod_{k \neq j} [p_k, \infty)$, such that
\\ for every seller $k \in [m]$, we have $R_k(\qvec)/R_{\all}(\qvec) = p_k/\beta$.
\end{quote}

As hinted in \cref{sec:overview}, we will prove that such a $\pvec$ exists using KKM.
To do so, for each seller $j$, we construct the set $C_j$ of price vectors
where every deviation to a higher price has a market-share-preserving response:
\[ C_j \defeq \bigg\{\pvec \in \beta\Delta_m:\; \forall q_j \in [p_j, \infty),\;
    \exists \qvec_{-j} \in \!\!\!\!\!\prod_{k \in [m] \setminus \{j\}} \!\!\! [p_k, \infty) \text{ such that }
    \forall k \in [m], \frac{R_k(\qvec)}{R_{\all}(\qvec)} = \frac{p_k}{\beta} \bigg\}. \]
A price vector belonging to $\bigcap_{j=1}^m C_j$ would be $\kappa$-RARE.
To show that this intersection is non-empty, we use KKM.
We can prove that the sets $(C_j)_{j \in [m]}$ satisfy the covering condition
by using an inner KKM, as outlined in \cref{sec:overview}.
It remains to show that each $C_j$ is closed. But $C_j$ need not be closed!
We give a counterexample in \cref{sec:unbounded-unclosed}.

We work around this obstacle by restricting the deviation and response sets to be compact.
In particular, we pick a large constant $L_k \gg B$ for each seller $k$,
and we require $q_k \le L_k$ for all $k \in [m]$.
With this change, the set
\begin{equation}
\label{eq:cover-sets}
C_j \defeq \bigg\{\pvec \in \beta\Delta_m:\; \forall q_j \in [p_j, L_j],\;
    \exists \qvec_{-j} \in \!\!\!\!\!\prod_{k \in [m] \setminus \{j\}} \!\!\! [p_k, L_k] \text{ such that }
    \forall k \in [m], \frac{R_k(\qvec)}{R_{\all}(\qvec)} = \frac{p_k}{\beta} \bigg\}
\end{equation}
becomes closed (we prove this in \cref{thm:rare-exists-continuous-2}).

It remains to show how to construct a response $\qvec_{-j}$ for deviations $q_j > L_j$.
If a seller $j$ deviates to a sufficiently large price, then the other sellers should be able to
easily out-compete her and drive down her revenue below $\nu p_j$.
But it's not obvious how to do this while maintaining individual rationality.
We do this by using another KKM!

\subsection{Market-Share-Restoring Responses Using KKM}
\label{sec:exist:kkm}

To apply KKM on the sets $(C_j)_{j \in [m]}$ as defined by \cref{eq:cover-sets},
we must prove the covering condition, i.e., for any $\pvec \in \beta\Delta_m$,
we have $\pvec \in C_j$ for some $j \in \support(\pvec)$.
Intuitively, this means that some seller $j$ with a positive price
should be unable to increase her market share.

We begin by proving a more general result: that some seller should be unable to \emph{gain}
by increasing prices, for an arbitrary definition of gain, not just `increase in market share'.
We prove this using KKM.
(This general result will be useful again later when analyzing responses to very high deviations,
i.e., the $q_j > L_j$ case.)

\begin{lemma}
\label{thm:general-inner-kkm}
Let $J \neq \emptyset$ be a finite set of sellers. Let $\ell_j < h_j$ for all $j \in J$.
Let $U$ be the union of the $|J|$ upper faces of the box $\prod_{j \in J} [\ell_j, h_j]$, i.e.,
\[ U \defeq \bigg\{\xvec \in \prod_{j \in J} [\ell_j, h_j]: x_i = h_i \text{ for some } i \in J\bigg\}. \]
For every $j \in J$, let $g_j: U \to \mathbb{R}$ be a continuous function
(the gain function) where for all $\xvec \in U$, the following hold:
\begin{itemize}
\item (Total gain is non-negative) $\sum_{j \in J} g_j(\xvec) \ge 0$.
\item (Non-positive gain at $\ell_j$) For all $j \in J$, $x_j = \ell_j \implies g_j(\xvec) \le 0$.
\end{itemize}
Then some $\xvec^* \in U$ satisfies $g_j(\xvec^*) \ge 0$ for all $j \in J$,
i.e., we get a non-negative gain for all $j \in J$.
\end{lemma}
\begin{proof}
For any $\theta \in \Delta_J \defeq \{\xvec \in \mathbb{R}_{\ge 0}^J: \sum_{j \in J} x_j = 1\}$,
shoot a ray from $\bm{\ell} = (\ell_j)_{j \in J}$ in the direction $\theta$,
and let $y(\theta)$ be the point on $U$ that the ray hits. Formally,
\[ y(\theta) \defeq \bm{\ell} + \tau(\theta)\theta, \quad \text{where} \quad
    \tau(\theta) \defeq \bigg(\max_{j \in J} \frac{\theta_j}{h_j - \ell_j}\bigg)^{-1}. \]
For all $j \in J$, define
\[ S_j \defeq \{\theta \in \Delta_J: g_j(y(\theta)) \ge 0 \}. \]
We now use KKM (\cref{thm:kkm}) on the simplex $\Delta_J$
to show that $\bigcap_{j \in J} S_j \neq \emptyset$.

\textbf{Closedness.}
For all $j \in J$, $y$ and $g_j$ are continuous, so $g_j \circ y$ is continuous.
Since $S_j$ is the superlevel set of $g_j \circ y$, it is closed
(see \cref{thm:superlevel-closed} in \cref{sec:exist-extra:continuous} for details).

\textbf{Covering.}
Suppose the covering condition doesn't hold,
i.e., $\exists \theta \in \Delta_J$ such that $\theta \not\in S_j$
for all $j \in \support(\theta) \defeq \{j \in J: \theta_j > 0\}$.
Hence, $\theta_j > 0 \implies g_j(y(\theta)) < 0$ for all $j \in J$.
For $j \in J \setminus \support(\theta)$, we have
$\theta_j = 0 \implies y(\theta)_j = \ell_j \implies g_j(y(\theta)) \le 0$.
Add all of these inequalities to get $\sum_{j \in J} g_j(y(\theta)) < 0$,
which is a contradiction. Thus, the covering condition holds.

By KKM, $\bigcap_{j \in J} S_j \neq \emptyset$. Let $\theta^* \in \bigcap_{j \in J} S_j$.
Let $\xvec^* \defeq y(\theta^*)$. Then $\xvec^* \in U$ and $g_j(\xvec^*) \ge 0$ for all $j \in J$.
\end{proof}

We now prove that on the upper boundary of the box $\prod_{j=1}^m [p_j, z_j]$,
one can find prices $\qvec$ that restore the market shares.
We do this by applying \cref{thm:general-inner-kkm} with
the gain function set to the increase in market share.

\begin{restatable}[market shares on box boundary]{lemma}{thmPropBox}
\label{thm:prop-box}
Let $A: \mathbb{R}_{\ge 0}^m \to [0, 1]^{n \times m}$ be a continuous allocation rule.
Let $\pvec \in \beta\Delta_m$ be a price vector
and $M^+ \defeq \{j \in [m]: p_j > 0\}$ be the positive-priced goods.

Let $\zvec \in \prod_{j \in M^+} (p_j, \infty) \times \{0\}^{[m] \setminus M^+}$
be a vector of upper bounds on increased prices.
Then there is another vector $\qvec$ on the boundary of the box $\prod_{j=1}^m [p_j, z_j]$
that preserves market shares, i.e., all of the following hold:
\begin{tightenum}
\item $p_j \le q_j \le z_j$ for all $j \in [m]$.
\item $q_i = z_i$ for some $i \in M^+$.
\item $R_{\all}(\qvec) \ge R_{\all}(\pvec) = \nu\beta$.
\item $\displaystyle \frac{R^A_j(\qvec)}{R_{\all}(\qvec)} = \frac{p_j}{\beta} = \frac{R^A_j(\pvec)}{R_{\all}(\pvec)}$
    for all $j \in [m]$.
\end{tightenum}
\end{restatable}
\begin{proof}
We get $R^A_j(\pvec) = \nu p_j$ by \cref{thm:accounting},
so $R_{\all}(\pvec) = \nu \sum_{j=1}^m p_j = \nu \beta$.

Let
\[ U \defeq \bigg\{\qvec \in \prod_{j \in M^+} [p_j, z_j] \times \{0\}^{[m] \setminus M^+}:
    q_i = z_i \text{ for some } i \in M^+\bigg\} \]
be the upper faces of the price-increase box,
and for each $j \in M^+$, let the gain function $g_j$ be the increase in market share, i.e.,
\[ g_j(\qvec) \defeq \frac{R_j(\qvec)}{R_{\all}(\qvec)} - \frac{p_j}{\beta}. \]
$R_{\all}$ is monotonic in prices (\cref{thm:rall-mono}),
so we get $R_{\all}(\qvec) \ge R_{\all}(\pvec) = \nu\beta$ for all $\qvec \in U$.
Then $g_j$ is continuous and $\sum_{j \in M^+} g_j(\qvec) = 0$.
Moreover, if $q_j = p_j$, then
\[ R_j(\qvec) \le \nu p_j = R_j(\pvec)
    \implies \frac{R_j(\qvec)}{R_{\all}(\qvec)} \le \frac{R_j(\pvec)}{R_{\all}(\pvec)}, \]
so $g_j(\qvec) \le 0$.

Hence, we can apply \cref{thm:general-inner-kkm}.
This gives us $\qvec \in U$ such that $g_j(\qvec) \ge 0$ for all $j \in M^+$.
Since $\sum_{j \in M^+} g_j(\qvec) = 0$, we have $g_j(\qvec) = 0$ for all $j \in M^+$.
Hence, $R_j(\qvec)/R_{\all}(\qvec) = p_j/\beta$ for all $j \in M^+$,
and this also holds for $j \in [m] \setminus M^+$ since $R_j(\qvec) = 0 = p_j$.
\end{proof}

We now apply KKM again to show that the sets $(C_j)_{j \in [m]}$ (as defined by \cref{eq:cover-sets})
have a common point $\pvec$. The covering condition for the KKM follows from \cref{thm:prop-box}.
At such a $\pvec$, for any seller $j \in [m]$ and deviation $q_j \in [p_j, L_j]$,
there is a market-share-restoring response.

\begin{lemma}
\label{thm:rare-exists-continuous-2}
Let $A: \mathbb{R}_{\ge 0}^m \to [0, 1]^{n \times m}$ be a continuous allocation rule
and $L_j \ge \beta$ for each seller $j \in [m]$.
%
There exist prices $\pvec \in \beta\Delta_m$ such that for every seller $j \in [m]$,
and every deviation $q_j \in [p_j, L_j]$ by seller $j$,
there exists a response $\qvec_{-j} \in \prod_{k \neq j} [p_k, L_k]$ such that
market shares are preserved, i.e.,
\[ \frac{R_k^A(\qvec)}{R_{\all}(\qvec)} = \frac{p_k}{\beta}
    \text{ for every seller } k \in [m]. \]
\end{lemma}
\begin{proof}
For every $j \in [m]$, let $C_j$ be the price vectors where every bounded deviation by seller $j$
can be met with a market-share-restoring response, i.e.,
\[ C_j \defeq \bigg\{\pvec \in \beta\Delta_m:\; \forall q_j \in [p_j, L_j],\;
    \exists \qvec_{-j} \in \!\!\!\!\!\prod_{k \in [m] \setminus \{j\}} \!\!\! [p_k, L_k] \text{ such that }
    \forall k \in [m], \frac{R_k(\qvec)}{R_{\all}(\qvec)} = \frac{p_k}{\beta} \bigg\}. \]
Our goal is to show that $\bigcap_{j=1}^m C_j \neq \emptyset$.
We prove this using KKM (\cref{thm:kkm}).

\textbf{Closedness.}
We first show that $C_j$ is closed. Let $\pvec^{(t)} \in C_j$ converge to $\pvec \in \beta\Delta_m$.
Fix any $q_j \in [p_j, L_j]$, and set $q_j^{(t)} \defeq \max(q_j, p_j^{(t)})$.
Since $\pvec^{(t)} \in C_j$, there is a response
$\qvec_{-j}^{(t)} \in \prod_{k \ne j} [p_k^{(t)}, L_k]$ such that
\[ \frac{R_k(\qvec^{(t)})}{R_{\all}(\qvec^{(t)})}
    = \frac{p_k^{(t)}}{\beta} \qquad \text{for every } k \in [m]. \]
The responses lie in the compact box $\prod_{k \ne j} [0, L_k]$,
so they must have a convergent subsequence.
Let $\qvec_{-j}$ be the limit of that subsequence.

By \cref{thm:rall-mono}, $R_{\all}(\qvec^{(t)}) \ge R_{\all}(\pvec^{(t)})$,
since $\qvec^{(t)} \ge \pvec^{(t)}$ coordinatewise.
Continuity of $A$ implies continuity of all revenue functions,
so taking limits gives us $R_{\all}(\qvec) \ge R_{\all}(\pvec) = \nu\beta$,
$q_k \ge p_k$ for all $k \in [m]$, and
\[ \frac{R_k(\qvec)}{R_{\all}(\qvec)} = \frac{p_k}{\beta}
    \qquad \text{for every }k \in [m]. \]
Thus, $\qvec_{-j} \in \prod_{k \neq j} [p_k, L_k]$ is a market-share-restoring response
to deviation $q_j$ and initial prices $\pvec$.
Since $q_j$ was arbitrary, we get $\pvec \in C_j$. Thus $C_j$ is closed.

\textbf{Covering.}
Fix $\pvec \in \beta\Delta_m$ and let $S \defeq \{j \in [m]: p_j > 0\}$ be its support.
Assume that the covering condition doesn't hold, i.e., $\pvec \not\in C_j$ for all $j \in S$.
For each $j \in S$, $\pvec \not\in C_j$ implies
\[ \exists q^{(j)}_j \in [p_j, L_j] \text{ such that }
    \forall \qvec^{(j)}_{-j} \in \!\!\!\!\!\prod_{k \in [m] \setminus \{j\}} \!\!\! [p_k, L_k],\;
    \exists k \in [m] \text{ such that } \frac{R_k(\qvec^{(j)})}{R_{\all}(\qvec^{(j)})} \neq \frac{p_k}{\beta}. \]
Let $\zvec \in \mathbb{R}^m$ be a vector where $z_j \defeq q^{(j)}_j$ for $j \in S$
and $z_j = 0$ for $j \in [m] \setminus S$.

Suppose $z_j = p_j$ for some $j \in S$. Then for $\qvec^{(j)}_{-j} = \pvec_{-j}$ and every $k \in [m]$, we get
\[ \frac{R_k(\qvec^{(j)})}{R_{\all}(\qvec^{(j)})} = \frac{R_k(\pvec)}{R_{\all}(\pvec)} = \frac{p_k}{\beta}, \]
which is a contradiction. Hence, $z_j > p_j$ for all $j \in S$.

By \cref{thm:prop-box}, there exist prices $\qvec$ such that
\begin{tightenum}
\item $p_j \le q_j \le z_j$ for all $j \in [m]$.
\item $q_i = z_i$ for some $i \in S$.
\item $R_{\all}(\qvec) \ge R_{\all}(\pvec) = \nu\beta$.
\item $\displaystyle \frac{R^A_j(\qvec)}{R_{\all}(\qvec)} = \frac{p_j}{\beta}$ for all $j \in [m]$.
\end{tightenum}
Picking $\qvec^{(i)}_{-i} = \qvec_{-i}$ gives us a contradiction.
Thus, the covering condition holds, so KKM applies.
\end{proof}

\subsection{The Futility of Large Deviations}
\label{sec:exist:overpricing}

\Cref{thm:rare-exists-continuous-2} gives us prices $\pvec$ such that
if any seller $j$ deviates to price $q_j \in [p_j, L_j]$,
there is a market-share-restoring response.
We now show that by defining $L_j$ appropriately,
we can construct a deterring response for $q_j \in (L_j, \infty)$.
To do this, we introduce additional parameters.

\paragraph{Security Parameters.}
The first of these are \emph{security parameters} $\sigma_j$ and $\rho_j$ for each seller $j$.
These say that if seller $j$ prices her good at $\sigma_j$, her revenue is at least $\rho_j$.
The following lemma proves that these parameters always exist.

\begin{restatable}{lemma}{thmSecurity}
\label{thm:security}
For any seller $j \in [m]$, there exist constants $\rho_j > 0$ and $\sigma_j \ge \rho_j/\nu$ such that
for any response $\qvec_{-j} \in \mathbb{R}_{\ge 0}^{[m]\setminus\{j\}}$ by the other sellers
and any allocation $\xvec \in \Dem((\sigma_j, \qvec_{-j}))$, we have
$R_j((\sigma_j, \qvec_{-j}), \xvec) \ge \rho_j$.
\end{restatable}
\begin{proof}[Proof sketch]
Intuitively, if seller $j$ sets a sufficiently small price for her good,
then every other seller falls into two categories:
\begin{tightenum}
\item \emph{cheap sellers}: these sellers also have a sufficiently low price,
    so they can extract only a small amount of revenue from the buyers,
    which leaves some budget for $j$.
\item \emph{expensive sellers}: these sellers charge a high price, so buyers prefer $j$ to these sellers.
\end{tightenum}
Thus, seller $j$ is guaranteed a minimum revenue.
We defer the full proof to \cref{sec:exist-extra:market-params}.
\end{proof}

\paragraph{Overpricing Parameter.}
Next, we study the \emph{overpricing parameter} $L_j$ for each seller $j$.
If the other sellers are trying to extract all budget $B$,
and if seller $j$ prices her good more than $L_j$, then she earns almost nothing.

\begin{restatable}{lemma}{thmOverpricing}
\label{thm:overpricing}
For any seller $j \in [m]$, define $\sigma_j$ and $\rho_j$ as in \cref{thm:security}.
Then there exists a constant $L_j \ge \max(B, \sigma_j)$ such that
for any prices $\qvec \in \mathbb{R}_{\ge 0}^m$ and allocation $\xvec \in \Dem(\qvec)$, we have
\[ q_j > L_j \text{ and } \max_{k \in [m] \setminus \{j\}} q_k = B
    \implies R_j(\qvec, \xvec) \le \frac{\rho_j}{2\kappa}. \]
\end{restatable}
\begin{proof}[Proof sketch]
Intuitively, if seller $j$'s price is too high, buyers will prefer the other sellers,
and since their prices are large enough to extract almost all buyer budget,
they will do so, leaving almost nothing for seller $j$.
We defer the full proof to \cref{sec:exist-extra:market-params}.
\end{proof}

The security parameters and the overpricing parameters together show
why it doesn't make sense to price too high:
seller $j$ can earn $\rho_j$ by pricing at $\sigma_j$,
but if she deviates to a price $q_j > L_j$,
then for some response $\qvec_{-j}$, she earns $\rho_j/(2\kappa) < \rho_j$.
It remains to show that such a response $\qvec_{-j}$ can be individually rational,
which we show in \cref{thm:large-dev-deter} using \cref{thm:general-inner-kkm},
which in turn uses KKM.

\begin{lemma}
\label{thm:large-dev-deter}
Let $A: \mathbb{R}_{\ge 0}^m \to [0, 1]^{n \times m}$ be a continuous allocation rule.
Define $(L_j)_{j \in [m]}$ as in \cref{thm:overpricing}.
%
Let $\pvec \in \beta\Delta_m$ be a price vector and $j \in [m]$ be a seller such that $p_j \ge \beta\rho_j/B$.
For every deviation $q_j > L_j$ by seller $j$, there exists a response
$\qvec_{-j} \in \prod_{k \neq j} [p_k, B]$ such that
\begin{tightenum}
\item ($j$ loses) $R^A_j(\qvec) < \nu p_j$.
\item (others gain) $R^A_k(\qvec) \ge \nu p_k$ for all $k \in [m] \setminus \{j\}$.
\end{tightenum}
\end{lemma}
\begin{proof}
Since $p_j > 0$, for all $k \in [m] \setminus \{j\}$, we get $p_k \le \beta - p_j < B$. Let
\[ U \defeq \bigg\{\qvec \in \{q_j\} \times \!\!\!\!\! \prod_{k \in [m] \setminus \{j\}} \!\!\! [p_k, B]:
    \max_{k \in [m] \setminus \{j\}} q_k = B\bigg\}. \]
For all $k \in [m] \setminus \{j\}$, define the gain function $g_k(\qvec) \defeq R_k(\qvec) - \nu p_k$.
Then $g_k$ is continuous.
If $q_k = p_k$ for some $k \in [m] \setminus \{j\}$, we have $R_k(\qvec) \le \nu p_k$, so $g_k(\qvec) \le 0$.

By \cref{thm:overpricing} (overpricing), for any $\qvec \in U$,
we get $R_j(\qvec) \le \rho_j/(2\kappa)$.
Also, $p_j \ge \beta\rho_j/B = \rho_j/(\kappa\nu)$.
Hence, $R_j(\qvec) \le \nu p_j / 2 < \nu p_j$.

Since $R_{\all}$ is monotonic in prices (\cref{thm:rall-mono}),
we get $R_{\all}(\qvec) \ge R_{\all}(\pvec) = \nu\beta$. Thus,
\[ \sum_{k \in [m] \setminus \{j\}} R_k(\qvec) = R_{\all}(\qvec) - R_j(\qvec)
    > \nu\beta - \nu p_j = \sum_{k \in [m] \setminus \{j\}} \nu p_k. \]
Hence, $\sum_{k \in [m] \setminus \{j\}} g_k(\qvec) > 0$.
%
Thus, we can apply \cref{thm:general-inner-kkm}.
This gives us $\qvec \in U$ such that $g_k(\qvec) \ge 0$ for all $k \in [m] \setminus \{j\}$,
i.e., $R_k(\qvec) \ge \nu p_k$.
\end{proof}

\subsection{RARE for Continuous Allocation Rules}
\label{sec:exist:continuous}

We are now ready to formally prove our main result:
the existence of a $\kappa$-RARE price vector
when the allocation rule is continuous.
In fact, we show something stronger:
\begin{tightenum}
\item We show that the sellers together extract a $1/\kappa$ fraction of the total buyer budget.
\item We show that the $\kappa$-RARE price vector $\pvec$ lies in the simplex $\beta\Delta_m$,
    and show a positive lower bound on each coordinate.
\item We show that the response $\qvec_{-j}$ to seller $j$'s deviation
    lies in the box $\prod_{k \neq j} [p_k, L_k]$.
\end{tightenum}
Note that we do not require the utilities to be concave.

\begin{theorem}
\label{thm:rare-exists-continuous}
Let $\Ical = ([n], [m], (u_i)_{i=1}^n, (b_i)_{i=1}^n)$ be a market instance
and $A: \mathbb{R}_{\ge 0}^m \to [0, 1]^{n \times m}$ be a continuous allocation rule.
Define parameters $(\rho_j)_{j=1}^m$ and $(L_j)_{j=1}^m$ as in \cref{sec:exist:overpricing}.
%
Then there exists a price vector $\pvec$ in the truncated simplex
$\{\yvec \in \beta\Delta_m: y_j \ge \beta\rho_j/B \text{ for all } j \in [m]\}$
such that for every seller $j \in [m]$ and every deviation $q_j \in \mathbb{R}_{\ge 0}$ by seller $j$,
there exists a response $\qvec_{-j} \in \prod_{k \in [m]\setminus\{j\}} [p_k, L_k]$ such that
\begin{tightenum}
\item (deviator $j$ doesn't gain much) $R_j^A(\qvec) \le \kappa \cdot \nu p_j = \kappa \cdot R_j^A(\pvec)$.
\item (responder $k$ doesn't lose) $R_k^A(\qvec) \ge \nu p_k \ge R_k^A((q_j, \pvec_{-j}))$
    for every seller $k \in [m] \setminus \{j\}$.
\end{tightenum}
\end{theorem}
\begin{proof}
Let $\pvec \in \beta\Delta_m$ be the prices obtained from
the double KKM argument (\cref{thm:rare-exists-continuous-2}).
Fix a seller $j \in [m]$. We first show that $p_j \ge \beta\rho_j/B$.

Suppose $p_j < \beta\rho_j/B$. Recall from \cref{thm:security} that $\rho_j/\nu \le \sigma_j$. Then
\[ p_j < \frac{\beta\rho_j}{B} = \frac{\rho_j}{\kappa\nu} \le \frac{\rho_j}{\nu} \le \sigma_j. \]
Let $q_j \defeq \sigma_j$. Since $q_j \in [p_j, L_j]$, by \cref{thm:rare-exists-continuous-2},
there exists a response $\qvec_{-j} \in \prod_{k \neq j} [p_k, L_k]$ such that
\[ R_j(\qvec) = \frac{p_j}{\beta} R_{\all}(\qvec) \le \frac{p_j}{\beta} B. \]
By \cref{thm:security} (security), we get $R_j(\qvec) \ge \rho_j$. Thus
\[ \rho_j \le \frac{p_j}{\beta} B \implies p_j \ge \frac{\beta\rho_j}{B}, \]
which is a contradiction. Hence, we have $p_j \ge \beta\rho_j/B$.

Now we show that there always exists a response $\qvec_{-j}$ satisfying the requirements of this theorem.
There are three cases, depending on the value of $q_j$.

\textbf{Case 1}: $q_j \in [p_j, L_j]$.
\\ Then by \cref{thm:rare-exists-continuous-2}, we get that there exists a response
$\qvec_{-j} \in \prod_{k \neq j} [p_k, L_k]$ such that for every $k \in [m]$, we have
\[ R_k(\qvec) = \frac{p_k}{\beta} R_{\all}(\qvec). \]
Thus,
\[ R_j(\qvec) = \frac{p_j}{\beta} R_{\all}(\qvec) \le \nu p_j \frac{B}{\nu\beta} = \kappa \nu p_j, \]
and for any $k \in [m] \setminus \{j\}$, using \cref{thm:rall-mono}, we have
\[ R_k(\qvec) = \frac{p_k}{\beta} R_{\all}(\qvec) \ge \frac{p_k}{\beta} R_{\all}(\pvec) = \nu p_k. \]
By \cref{thm:accounting}, we get $R_j(\pvec) = \nu p_j$ and $R_k((q_j, \pvec_{-j})) \le \nu p_k$.
Thus, $\qvec_{-j}$ satisfies the requirements of this theorem.

\textbf{Case 2}: $q_j < p_j$.
\\ Let $\qvec_{-j} \defeq \pvec_{-j}$. By \cref{thm:accounting}, we get
$R_j(\qvec) = \nu q_j < \nu p_j$ and for all $k \in [m] \setminus \{j\}$,
we get $R_k(\qvec) = \nu q_k = \nu p_k = R_k((q_j, \pvec_{-j}))$.
Thus, $\qvec_{-j}$ satisfies the requirements of this theorem.

\textbf{Case 3}: $q_j > L_j$.
\\ \Cref{thm:large-dev-deter} gives us a response $\qvec_{-j} \in \prod_{k \neq j} [p_k, B]$
such that $R_j(\qvec) \le \nu p_j$ and $R_k(\qvec) \ge \nu p_k$ for all $k \in [m] \setminus \{j\}$.
By \cref{thm:accounting}, we get $R_j(\pvec) = \nu p_j$ and $R_k((q_j, \pvec_{-j})) \le \nu p_k$.
Thus, $\qvec_{-j}$ satisfies the requirements of this theorem.
\end{proof}

\Cref{thm:rare-exists-continuous} implies the following simple corollary.

\begin{corollary}
\label{thm:rare-exists-continuous-corr}
For any market $\Ical$ and continuous allocation rule $A$,
there exists a $\kappa$-RARE price vector $\pvec$
at which the sellers' total revenue is $R_{\all}(\pvec) = B/\kappa$,
where $\kappa = 1$ in the rivalrous setting and
$\kappa = b_{\avg}/b_{\min}$ in the non-rivalrous setting.
\end{corollary}
\begin{proof}
\Cref{thm:rare-exists-continuous} gives us the $\kappa$-RARE price vector $\pvec \in \beta\Delta_m$.
By \cref{thm:accounting}, we get $R_{\all}(\pvec) = \nu\sum_{j=1}^m p_j = \nu\beta = B/\kappa$.
\end{proof}

\subsection{RARE for Non-Rivalrous Goods}
\label{sec:exist:non-rival}

In non-rivalrous markets, when the utility functions are strictly concave,
the demand set is a singleton, so there is a unique allocation rule.
Moreover, one can prove that such an allocation rule is continuous.
We defer the full proof to \cref{sec:exist-extra:non-rival}.

\begin{restatable}{theorem}{thmUniqueNrDemand}
\label{thm:unique-nr-demand}
Let $\Ical = ([n], [m], (u_i)_{i=1}^n, (b_i)_{i=1}^n)$ be a market instance.
Let $u_i$ be strictly concave for all $i \in [n]$.
Then for all prices $\pvec \in \mathbb{R}_{\ge 0}^m$,
the non-rivalrous demand $\DemNR(\pvec)$ is a singleton.
Moreover, the unique allocation rule $A$ is continuous.
\end{restatable}

Applying \cref{thm:rare-exists-continuous-corr} then shows that for strictly
concave utilities, a $(b_{\avg}/b_{\min})$-RARE always exists.
On the other hand, for linear utility functions,
there are instances where every allocation rule is discontinuous.

\begin{example}
\label{ex:lin-util-dem-discont}
Let there be two goods,
and a single buyer with budget 1 and utility function $u_1(\xvec) = x_1 + x_2$.
Let $\pvec \in (0, 1)^2$ be a price vector such that $p_1 + p_2 > 1$.
If $p_1 < p_2$, then $\Dem_1(\pvec) = \{(1, \frac{1-p_1}{p_2})\}$.
If $p_1 > p_2$, then $\Dem_1(\pvec) = \{(\frac{1-p_2}{p_1}, 1)\}$.
Thus, for any $\eps \in (0, 1/3)$, we have $\Dem_1((2/3-\eps, 2/3+\eps)) = \{(1, (1+3\eps)/(2+3\eps))\}$
and $\Dem_1((2/3+\eps, 2/3-\eps)) = \{((1+3\eps)/(2+3\eps), 1)\}$.
Thus, $\Dem_1$ is discontinuous.
\end{example}

However, we can still show that approximate RARE exists by adding
an infinitesimally small concave perturbation to the utilities.
We use $\vecOne$ to denote the allocation where every good is allocated to every buyer.

\begin{theorem}[\texorpdfstring{$(b_{\avg}/b_{\min})$}{(bavg/bmin)}-RARE for non-rivalrous markets]
\label{thm:nonrival-rare-exists}
Let $\Ical = ([n], [m], (u_i)_{i=1}^n, (b_i)_{i=1}^n)$ be a non-rivalrous market instance.
Suppose each utility $u_i: [0,1]^m \to \mathbb{R}_{\ge 0}$ is concave.
Let $\kappa \defeq b_{\avg}/b_{\min}$.
%
Then there exist prices $\pvec \in b_{\min}\Delta_m^{\circ}$ such that $\vecOne \in \DemNR(\pvec)$ and
for every seller $j \in [m]$ and deviation price $q_j \in \mathbb{R}_{\ge 0}$,
there exists a response $\qvec_{-j} \in \mathbb{R}_{\ge 0}^{[m]\setminus\{j\}}$ and allocation $\xvec \in \DemNR(\qvec)$
(where $\qvec \defeq (q_j, \qvec_{-j})$) such that:
\begin{tightenum}
\item Seller $j$'s revenue after the response is at most $\kappa$ times her original revenue, i.e.,
    $R_j(\qvec, \xvec) \le \kappa \cdot R_j(\pvec, \vecOne) = \kappa \cdot n p_j$.
\item No responder loses by responding, i.e., $R_k(\qvec, \xvec) \ge n p_k$
    for every responder $k \in [m] \setminus \{j\}$.
\end{tightenum}
\end{theorem}
\begin{proof}
For every $\eps \in [0, 1]$, define the perturbed instance
$\Ical^{(\eps)} \defeq ([n], [m], (u_i^{(\eps)})_{i=1}^n, (b_i)_{i=1}^n)$, where
\[ u_i^{(\eps)}(\xvec) \defeq u_i(\xvec) + \eps \sum_{k=1}^m x_k(2-x_k). \]
One can show that the parameters $\rho_j > 0$, $\sigma_j \ge \rho_j/\nu$,
and $L_j \ge \max(B, \sigma_j)$ can be picked to be independent of the perturbation.
See \cref{thm:robust-params} in \cref{sec:exist-extra:market-params}.

For $\eps_t = 1/t$, for each $t \in \N$ and each $i \in [n]$,
the utility function $u_i^{(\eps_t)}$ is strictly concave.
Hence, by \cref{thm:unique-nr-demand}, there is a unique continuous allocation rule
$A^{(\eps_t)}$ for each instance $\Ical^{(\eps_t)}$.
Apply \cref{thm:rare-exists-continuous} to obtain $\pvec^{(t)} \in b_{\min}\Delta_m$
and $p^{(t)}_j \ge b_{\min}\rho_j/B$.

By compactness, pass to a subsequence along which $\pvec^{(t)} \to \pvec \in b_{\min}\Delta_m$.
Relabel this subsequence and its corresponding perturbations using the same notation;
we still have $\eps_t \to 0$.
Fix this price subsequence before choosing a seller or deviation.
By the uniform lower bound, we get $p_j \ge b_{\min}\rho_j/B > 0$ for all $j \in [m]$.
Hence, every buyer can buy all the goods, so $\vecOne \in \DemNR(\pvec; \Ical)$.

Fix a seller $j \in [m]$ and deviation $q_j$. By \cref{thm:rare-exists-continuous},
we get a response $\qvec_{-j}^{(t)}$ in each instance $\Ical^{(\eps_t)}$ such that
$q^{(t)}_k \in [p_k^{(t)}, L_k]$ for all $k \in [m] \setminus \{j\}$.
Moreover, let $\xvec^{(t)} \defeq A^{(\eps_t)}(\qvec^{(t)})$ be the
unique element in $\DemNR(\qvec^{(t)}; \Ical^{(\eps_t)})$.
\Cref{thm:rare-exists-continuous} also guarantees
$R_j(\qvec^{(t)}, \xvec^{(t)}) \le \kappa n p_j^{(t)}$
and $R_k(\qvec^{(t)}, \xvec^{(t)}) \ge n p_k^{(t)}$ for all $k \in [m] \setminus \{j\}$.

The pairs $(\qvec_{-j}^{(t)}, \xvec^{(t)})$ lie in the compact set
\[ \bigg(\prod_{k \neq j} [0, L_k]\bigg) \times [0, 1]^{n \times m}. \]
Thus, there is a further subsequence along which
$(\qvec_{-j}^{(t)}, \xvec^{(t)}) \to (\qvec_{-j}, \xvec)$.
Along this same subsequence, $\pvec^{(t)} \to \pvec$ and $\eps_t \to 0$ still hold.
This further subsequence may depend on $j$ and $q_j$.
Let $\qvec \defeq (q_j, \qvec_{-j})$; since the deviation price is fixed, $\qvec^{(t)} \to \qvec$.
Then $\qvec_{-j} \in \prod_{k \neq j} [p_k, L_k]$, and by \cref{thm:demand-limit}, we get $\xvec \in \DemNR(\qvec; \Ical)$.
By the continuity of $R_j(\cdot, \cdot)$, we get $R_j(\qvec, \xvec) \le \kappa n p_j$
and $R_k(\qvec, \xvec) \ge n p_k$ for all $k \in [m] \setminus \{j\}$.
\end{proof}

\subsection{RARE for Rivalrous Goods}
\label{sec:exist:rival}

For non-rivalrous goods, strictly concave utilities determine a unique allocation rule.
For merely concave utilities, \cref{thm:nonrival-rare-exists} gives a $\kappa$-RARE guarantee
without a fixed allocation rule: a demand allocation can be chosen for each response.

For rivalrous goods, the demand set can contain multiple allocations even with strictly concave utilities.
E.g., if there are two buyers and two goods, priced such that no buyer can afford either good completely,
then the demand set $\DemR(\pvec)$ can have very different allocations.
We restrict the choice of allocation by maximizing a fixed \emph{welfare function}.
When its maximizer is unique at every price vector, this determines an allocation rule.

\begin{definition}[\texorpdfstring{$W$}{W}-rivalrous demand]
Let $W: \mathbb{R}_{\ge 0}^n \to \mathbb{R}$ be a continuous
and strictly monotone function, which we call the \emph{welfare function}.
Define the $W$-rivalrous demand at prices $\pvec \in \mathbb{R}_{\ge 0}^m$ as
\[ \Dem^W(\pvec) \defeq \argmax_{\xvec \in \AffR(\pvec)} W(u_1(\xvec_1), \ldots, u_n(\xvec_n)). \]
(The argmax is well-defined: since $W$ is continuous and $\AffR(\pvec)$ is a compact set, a maximum always exists.)
\end{definition}

Recall that $\AffR(\pvec)$ is the set of allocations where each buyer can afford her bundle.
Since $W$ is strictly monotone, the maximizer allocation $\xvec$ is buyer-PO, so $\Dem^W(\pvec) \subseteq \DemR(\pvec)$.

\Cref{ex:lin-util-dem-discont} shows that in a market with two sellers and one buyer
with linear utilities, every allocation rule can be discontinuous.
Thus, we study strictly concave utility functions.
For such utilities, when $W$ is concave, then $\Dem^W(\pvec)$ is a singleton
and the corresponding unique allocation rule is continuous.
We defer the full proof to \cref{sec:exist-extra:rival}.

\begin{restatable}{theorem}{thmUniqueRivalDemand}
\label{thm:unique-rival-demand}
Let $\Ical = ([n], [m], (u_i)_{i=1}^n, (b_i)_{i=1}^n)$ be a market instance.
Let $u_i$ be strictly concave for all $i \in [n]$.
Let $W: \mathbb{R}_{\ge 0}^n \to \mathbb{R}$ be a continuous, strictly monotone, concave function.
Then for all prices $\pvec \in \mathbb{R}_{\ge 0}^m$,
the $W$-rivalrous demand $\Dem^W(\pvec)$ is a singleton.
Moreover, the unique allocation rule $A$ is continuous.
\end{restatable}

We give several examples of concave welfare functions:
\begin{enumerate}
\item \textbf{Utilitarian social welfare}: Set $W(\xvec) = \sum_{i=1}^n x_i$.
\item \textbf{Weighted social welfare}: Set $W(\xvec) \defeq \wvec^T\xvec$, where $\wvec \in \mathbb{R}_{>0}^n$.
\item \textbf{Priority order}:
    Set $W(\xvec) \defeq \wvec^T\xvec$, where $w_i \defeq M^{-i}$ for each buyer $i \in [n]$
    and $M$ is a very large number. This approximates the requirement that we give the maximum
    possible utility to buyer 1, then the maximum possible utility to buyer 2, and so on.
    We can renumber the buyers \wLoG{} to accommodate any priority order.
\item \textbf{Weighted Nash Social Welfare}: For a very small $\eps > 0$ and $\wvec \in \Delta_n^{\circ}$,
    \[ W(\xvec) \defeq \prod_{i=1}^n x_i^{w_i} + \eps \sum_{i=1}^n x_i, \]
    gives an approximation to the Weighted Nash Social Welfare.
    Adding $\eps$ times the utilitarian social welfare ensures that $W$ is strictly monotone.
\item \textbf{Leximin Welfare}: For $0 < \eps \ll 1$,
    \[ W(\xvec) \defeq \sum_{i=1}^n \eps^{i-1} x_{[i]}, \]
    where $x_{[i]}$ is the $i\Th$ smallest coordinate in $\xvec$,
    gives an approximation to the Leximin Welfare.
\end{enumerate}
For all of the above welfare functions, there is a corresponding unique allocation rule
(when utilities are strictly concave) which is continuous.
For that allocation rule, a RARE always exists by \cref{thm:rare-exists-continuous-corr}.

When utilities are merely concave, $\Dem^W(\pvec)$ may contain several allocations,
so fixing $W$ need not determine a unique allocation rule.
The following theorem gives a version of RARE without a fixed allocation rule.
It fixes $W$ and the initial allocation $\xvec$, and allows a welfare-maximizing allocation $\zvec$
to be chosen separately for each seller's deviation and response.
We obtain this guarantee by perturbing the utilities to make them strictly concave and taking a limit.
The proof is similar to \cref{thm:nonrival-rare-exists},
so we defer the full proof to \cref{sec:exist-extra:rival}.

\begin{restatable}{theorem}{thmRivalRareExists}
\label{thm:rival-rare-exists}
Let $\Ical = ([n], [m], (u_i)_{i=1}^n, (b_i)_{i=1}^n)$ be a market instance.
Suppose each utility $u_i: [0,1]^m \to \mathbb{R}_{\ge 0}$ is concave.
Fix a continuous, strictly monotone, concave function $W: \mathbb{R}_{\ge 0}^n \to \mathbb{R}$.
Then there exist prices $\pvec \in B\Delta_m^{\circ}$ and an allocation $\xvec \in \Dem^W(\pvec)$
such that for every seller $j \in [m]$ and deviation price $q_j \in \mathbb{R}_{\ge 0}$,
there exists a response $\qvec_{-j} \in \mathbb{R}_{\ge 0}^{[m]\setminus\{j\}}$ and allocation $\zvec \in \Dem^W(\qvec)$
(where $\qvec \defeq (q_j, \qvec_{-j})$) such that:
\begin{tightenum}
\item $j$'s revenue after the response is at most her original revenue, i.e.,
    $R_j(\qvec, \zvec) \le R_j(\pvec, \xvec) = p_j$.
\item No responder loses by responding, i.e., $R_k(\qvec, \zvec) \ge p_k$
    for every responder $k \in [m] \setminus \{j\}$.
\end{tightenum}
\end{restatable}
